\documentclass[12pt]{article}
\usepackage{geometry}
\usepackage{amsmath}
\usepackage[utf8]{inputenc}
\usepackage{xcolor}
\usepackage{hyperref}
\usepackage{parskip}
\usepackage[hang, flushmargin, bottom]{footmisc}

\geometry{verbose,tmargin=1in,bmargin=1in,lmargin=1in,rmargin=1in,nomarginpar}
\hypersetup{breaklinks=true,hypertexnames=false,colorlinks=true,citecolor=teal,urlcolor=teal}
\setlength{\parskip}{0.9em}

\begin{document}

\noindent Lucas Schmitz\footnote{Yale University. \texttt{lucas.schmitz@yale.edu}}\\
PhD Candidate in Economics, Yale University\\
February 22, 2026

\bigskip

\noindent Gerencia de Personas\\
Banco Central de Chile

\bigskip

\noindent Dear Selection Committee,

I am a PhD candidate in Economics at Yale University and I am writing to apply to the Programa de Visitas de Trabajo 2026. My dissertation studies competition in consumer credit markets, with a focus on how information asymmetries and switching costs shape lender behavior, credit pricing, and borrower welfare. I would like to conduct an in-person visit at the Banco Central between June 20 and July 31, 2026.

\noindent\textbf{Area of interest and proposed activity}

The proposed visit would center on estimating the effects of Law 21.680, which established Chile's Consolidated Debt Registry (REDEC) in July 2024. Prior to REDEC, only banks and the largest savings cooperatives (CACs) were required to report to---and permitted to access---the CMF's N\'{o}mina de Deudores, while non-bank credit providers (OCNBs)---including retail lenders, cajas de compensaci\'{o}n, and factoring and leasing firms, which account for over 60\% of consumer credit placements---were excluded from the registry. REDEC extends both the reporting obligation and registry access to all credit providers, changing the information available to lenders when evaluating borrowers. This creates a clean before-and-after comparison to study how broader information sharing affects credit pricing, competition between bank and non-bank lenders, and household over-indebtedness.

The deliverable for the visit would be an internal research seminar presenting preliminary empirical findings from this project. The analysis would draw on the administrative credit records held at the Banco Central, using structural Industrial Organization methods to estimate the competitive and distributional effects of the reform across borrower segments.

\noindent\textbf{Motivation}

The Banco Central is uniquely positioned to host this project. The administrative data required to study REDEC rigorously---covering borrower-level debt positions across all regulated lenders---is held at the institution, and access to it is a prerequisite for credible empirical work. Beyond data access, the Bank has followed this policy question for a long time: the Informes de Estabilidad Financiera for the second semester of 2018 and the first semester of 2019 documented the financial stability risks arising from Chile's fragmented debt information system, and researchers at the Bank have been central to the literature on consumer credit markets in Chile. A visit would allow me to work alongside that expertise directly, sharpen the empirical design with input from staff familiar with the data, and produce findings of direct relevance to the Bank's financial stability and research mandates.

\bigskip

\noindent I look forward to the possibility of contributing to the work of the institution. Please do not hesitate to contact me if you need any additional information.

\bigskip

\noindent Sincerely,

\bigskip
\bigskip

\noindent Lucas Schmitz

\end{document}